\appendix

\section{Proof details}
\label{app:proofs}

\subsection{Proof of Lemmas \ref{lem:cournot-continuation} and \ref{lem:global-concavity}}

For fixed marginal costs, the first derivative of firm \(j\)'s quantity payoff is
\[
\frac{\partial\pi_j}{\partial q_j}
=
a-c_j-bQ-bq_j.
\]
The potential \eqref{eq:potential} has the same gradient and Hessian
\[
\nabla^2\Phi=-b(I+\mathbf 1\mathbf 1^\top),
\]
which is negative definite for \(b>0\).  Maximization of \(\Phi\) over the convex set \(\mathbb R_+^N\) therefore has a unique KKT solution.  The KKT conditions yield \eqref{eq:pS}--\eqref{eq:active-set-q}, proving Lemma~\ref{lem:cournot-continuation}.

Now fix rivals' sourcing and vary \(x=i_j\).  On a branch with \(m\) active firms including \(j\), \(z_j=p-c_j\) satisfies
\[
\frac{dz_j}{dx}=H\frac{m}{m+1}
\]
and operating profit is \(z_j^2/b\).  Hence \eqref{eq:branch-curvature}.  Since
\[
\frac{m}{m+1}\le\frac{N}{N+1},
\]
condition \eqref{eq:source-soc} makes every smooth branch strictly concave.  At own entry \(z_j=0\), so the derivative is continuous.  At a rival-exit threshold, the derivative falls by
\[
\frac{2Hz_j}{bm(m+1)}>0.
\]
Thus the derivative is strictly decreasing globally, proving Lemma~\ref{lem:global-concavity}.

\subsection{Proof of Proposition \ref{prop:symmetric-cournot}}

At a symmetric profile \(s\), each firm produces
\[
q(s)=\frac{D+Hs}{b(N+1)}.
\]
The derivative of reduced payoff with respect to a unilateral sourcing choice, evaluated at symmetry, is
\[
\frac{2HN(D+Hs)}{b(N+1)^2}-2s.
\]
Its unconstrained zero is \eqref{eq:iCbar}.  Lemma~\ref{lem:global-concavity} makes this condition globally sufficient, while \(D>0\) rules out the lower bound.  Projection onto \([0,\phi]\) gives \eqref{eq:iCstar}.  Differentiating the continuous extension gives \eqref{eq:dIdN}.  Since \(N\) is economically integer-valued, the theorem states the resulting monotonic ordering only across admissible integer values satisfying the maintained source restrictions.

\subsection{Derivation of the source-duopoly response}

Normalize as in \eqref{eq:normalized}.  Against rival action \(y\), the both-active stationary point, rival-exit threshold, and monopoly stationary point are \eqref{eq:BR-components}.  Their relevant joins solve
\[
A(y_A)=U(y_A),
\qquad
M=U(y_M),
\]
which gives \eqref{eq:yA-yM}.  Global own-payoff concavity implies that the global unconstrained maximizer moves through the branch sequence reported in \eqref{eq:R}; projection onto the primitive sourcing interval gives \eqref{eq:Bphi}.  On the exit branch \(U'(y)=2\), which proves Proposition~\ref{prop:prop5}.

\subsection{Proof of Proposition \ref{prop:correspondence}}

A pure Stage-I equilibrium solves
\[
x=B_\phi(y),\qquad y=B_\phi(x).
\]
If \(\rho>2/3\), the active branch has slope magnitude \(2/(9\rho-4)<1\), while cap and zero pieces have slope zero.  Hence
\[
T(x,y)=(B_\phi(y),B_\phi(x))
\]
is a contraction in the sup norm; its unique fixed point is symmetric and equals \((\min\{\phi,s\},\min\{\phi,s\})\).

At \(\rho=2/3\),
\[
B_\phi(y)=\min\{\phi,\max\{0,\delta-y\}\}.
\]
If \(\phi<\delta/2\), both best responses bind at the cap and the unique fixed point is \((\phi,\phi)\).  If \(\phi\ge\delta/2\), mutual best response is equivalent to \(x+y=\delta\) with both actions in \([0,\phi]\), giving \eqref{eq:knife-continuum}.

For \(4/9<\rho<2/3\), the symmetric fixed point is \(s\).  If \(\phi\le s\), the cap yields the unique fixed point \((\phi,\phi)\).  If \(\phi>s\), intersecting each flat, exit, active, zero, and cap segment of the complete graph with its transpose yields the symmetric point plus one ordered asymmetric pair.  Writing \(x_H=\min\{\phi,R(0)\}\) and \(x_L=R(x_H)\) gives \eqref{eq:three-eq}; no other branch intersection satisfies mutual best response.

\subsection{Proof of Proposition \ref{prop:literal-hotelling}}

Firm A's pure price payoff is
\[
n(p_A-c_A)s_A(p_A,p_B),
\]
with \(s_A\) given by \eqref{eq:clipped-share}.  The global best response consists of a zero-demand region, the interior quadratic maximizer, and the full-market boundary.  Solving mutual best responses gives the standard unique interior prices when \(|d|<3\tau\), the unique boundary point at equality, and \eqref{eq:literal-corner} when \(d>3\tau\).  In the latter case the high-cost firm receives zero demand for every \(z\in[c_A+3\tau,c_B]\); the low-cost firm's best response is \(z-\tau\).  The opposite cost ordering is symmetric.

\subsection{Proof of Proposition \ref{prop:no-loss}}

Under \eqref{eq:no-loss}, the corner continuation for \(d>3\tau\) is uniquely
\[
p_B=c_B,\qquad p_A=c_B-\tau,
\]
with the mirror image for \(d<-3\tau\).  Substitution into Stage-I payoffs gives \eqref{eq:no-loss-payoff}.  Against the symmetric interior candidate \(i_0=Hn/6\), the only potentially profitable nonlocal deviation is to the low-cost full-market branch.  Its gain is
\begin{equation}
\Delta(x)
=
-\frac{
36x^2-36Hnx+5H^2n^2+54n\tau
}{36}.
\label{eq:delta-x}
\end{equation}
The unconstrained maximizer is \(x_M=Hn/2\), with
\[
\Delta(x_M)
=
\frac{n(2H^2n-27\tau)}{18}.
\]
Thus a positive-gain region requires \(H^2n>27\tau/2\).  Combined with the source SOC \(H^2n<18\tau\), the lower zero of \(\Delta\) is exactly \eqref{eq:xminus}; in this region it lies inside the correct full-market branch.  A feasible strictly profitable deviation exists if and only if \(\phi>x_-\).  Equality produces a tie.  This proves \eqref{eq:no-loss-region}.  Substituting \(H=n=1\), \(\tau=1/15\), and \(x=1/2\) into \eqref{eq:delta-x} yields the exact gain \(1/90\).
